\documentclass[11pt]{article}
\usepackage[a4paper,margin=1in]{geometry}
\usepackage{amsmath,amssymb,mathtools,bm}
\usepackage{enumitem,booktabs,array}
\usepackage{tcolorbox}
\usepackage{hyperref}
\usepackage[protrusion=true,expansion=false]{microtype}
\hypersetup{colorlinks=true,linkcolor=blue,urlcolor=blue,citecolor=blue}
\setlist[itemize]{topsep=2pt,itemsep=2pt}
\setlist[enumerate]{topsep=2pt,itemsep=2pt}
\newtcolorbox{keybox}{colback=blue!3!white,colframe=blue!40!black,boxrule=0.4pt,arc=1mm}
\newtcolorbox{alertbox}{colback=red!3!white,colframe=red!40!black,boxrule=0.4pt,arc=1mm}
\newtcolorbox{answerbox}{colback=green!3!white,colframe=green!40!black,boxrule=0.4pt,arc=1mm}
\newtcolorbox{drillbox}{colback=yellow!5!white,colframe=orange!60!black,boxrule=0.4pt,arc=1mm}
\newcommand{\EV}{\mathbb{E}}
\newcommand{\Var}{\mathrm{Var}}
\newcommand{\Cov}{\mathrm{Cov}}
\newcommand{\blank}[1]{\underline{\hspace{#1}}}

\title{W09-02: Bakke, Mahmudi, Fernando \& Salas (2016, JFE) \\
\large ``The Causal Effect of Option Pay on Corporate Risk Management'' --- Active Recall Workbook}
\author{Banking \& Risk Management --- Exam Prep}
\date{\today}
\begin{document}\maketitle

\begin{keybox}
\textbf{Paper in one sentence.} Using FAS 123R (2006) --- which forced firms to expense options and caused heavy-option-users to cut option pay --- as an exogenous shock to CEO compensation convexity, BMFS show that firms whose CEOs lost option pay \emph{increased hedging} by about $\sim 20\%$, confirming Smith--Stulz (1985): convex CEO pay reduces hedging incentives, so cutting option pay \emph{raises} risk management.

\textbf{Placement.} Bridges Smith--Stulz theory (1985), GMS 1997 empirical determinants (W07-01), and the Gormley-Matsa-Milbourn (W09-01) natural-experiment tradition. Canonical causal evidence that comp-convexity causes hedging choices.
\end{keybox}

\section*{Round 1: Complete Walk-Through}

\subsection*{1.1 The accounting shock}
Prior to FAS 123R, firms did not expense stock options on income statements. After FAS 123R (effective 2006), options hit earnings immediately. Firms that had relied heavily on option pay faced a sudden reported-EPS cost, and many cut option grants sharply.

\subsection*{1.2 Identification}
DiD: firms with ex-ante high option-pay reliance (treated) vs.\ low option-pay reliance (control), pre vs.\ post FAS 123R:
\[
\text{Hedging}_{it}=\alpha_i+\delta_t+\beta\cdot(\text{HighPreOptPay}_i\cdot\text{Post}_t)+X_{it}'\gamma+\varepsilon_{it}.
\]
$\beta$ captures the change in hedging at firms whose CEOs lost option pay.

\subsection*{1.3 Data}
ExecuComp for CEO pay structure; FAS 10-K, 10-Q disclosures for hedging; Compustat for firm fundamentals. Sample 2000--2010, covering 3--4 years pre and post the shock.

\subsection*{1.4 Headline results}
\begin{itemize}
\item Treated firms reduce CEO option pay by $\sim$30--40\% post-shock.
\item Hedging (indicator and intensity) rises by $\sim 20\%$ at these firms relative to controls.
\item Effect is concentrated in firms with commodity, FX, and interest-rate exposures where hedging is most tractable.
\item Placebos and pre-trends support the causal interpretation.
\end{itemize}

\subsection*{1.5 Theoretical frame}
Smith--Stulz (1985): option-heavy pay makes CEO's utility \emph{convex} in firm value $\Rightarrow$ CEO prefers volatility $\Rightarrow$ hedges less. Cutting option pay \emph{reduces} convexity $\Rightarrow$ CEO hedges more. BMFS is the causal empirical translation of that prediction.

\subsection*{1.6 Caveats}
\begin{alertbox}
\begin{itemize}
\item FAS 123R affects more than just option pay (accounting, market perception); attribution is careful.
\item Hedging disclosures are sometimes qualitative; BMFS use multiple measures.
\item Option cut could co-move with other governance changes; robustness tests included.
\end{itemize}
\end{alertbox}

\section*{Round 2: Level 1 Fill-in}
\begin{enumerate}
\item Shock: FAS \blank{1cm}R (year \blank{1cm}) option expensing.
\item DiD treatment: high ex-ante \blank{1cm}-pay reliance.
\item Hedging change: $\sim$\blank{1cm}\% \blank{1cm}.
\item Theoretical prediction: Smith--Stulz \blank{1cm}.
\item Option-pay cut magnitude: $\sim$\blank{1cm}--\blank{1cm}\%.
\end{enumerate}

\section*{Round 3: Level 2 Prompts}
\begin{enumerate}
\item Explain why option pay makes CEO utility convex and how this affects hedging.
\item Write the DiD specification and the identifying assumption.
\item Why does FAS 123R produce an arguably exogenous shock to option pay?
\item Contrast BMFS's causal identification with GMS (1997) cross-section.
\end{enumerate}

\section*{Round 4: Minimal Scaffolding}
From a blank page: (i) shock, (ii) DiD, (iii) magnitudes, (iv) theoretical mapping, (v) two caveats.

\section*{Round 5: Blank Page}
Essay: does CEO pay design cause hedging? BMFS's causal evidence and its place in the literature.

\vspace{6pt}
\begin{drillbox}
\textbf{Speed Drill.} (1) Shock. (2) Treatment intensity. (3) Hedging direction/magnitude. (4) Theoretical paper the result maps to. (5) Data source for comp.
\end{drillbox}

\section*{Quick Reference \& Answer Key}
\begin{answerbox}
\begin{itemize}
\item \textbf{Shock.} FAS 123R, 2006 option expensing.
\item \textbf{Design.} DiD high vs.\ low ex-ante option pay.
\item \textbf{Effect.} Hedging $+\sim\!20\%$ at treated firms.
\item \textbf{Theory.} Smith--Stulz (1985): less convex pay $\Rightarrow$ more hedging.
\end{itemize}
\end{answerbox}

\begin{keybox}
\textbf{30-second exam pitch.} Bakke, Mahmudi, Fernando \& Salas (2016) causally identify the effect of CEO compensation convexity on corporate hedging. FAS 123R (2006) forced firms to expense options on the income statement, causing heavy-option-users to cut option pay by 30--40\% and thereby lowering CEO compensation convexity. DiD comparing high- vs.\ low-pre-shock-option firms shows that treated firms \emph{increased} hedging by roughly 20\% post-shock, with larger effects in firms facing commodity, FX or rate exposures. The result is the textbook confirmation of Smith--Stulz (1985): when CEO utility is convex in firm value, CEOs prefer volatility and hedge less; reducing that convexity flips the incentive, and managers hedge more. The paper complements Gormley--Matsa--Milbourn (2013) --- both show that compensation structure is the active adjustment margin for corporate risk.
\end{keybox}

\end{document}
